% =====================================================================
%  Small attention-head sets that carry simple tasks: study design.
%  Written as a paper; every result is a placeholder (\pending) until the runs
%  are made. Compile with Tectonic (XeTeX).
% =====================================================================
\documentclass[10pt,a4paper]{article}

\usepackage{fontspec}
\usepackage[margin=2.2cm]{geometry}
\usepackage{amsmath,amssymb}
\usepackage{booktabs,tabularx,array}
\usepackage{graphicx}
\usepackage{float}
\usepackage{enumitem}
\usepackage{xcolor}
\usepackage{caption}
\usepackage[round,authoryear]{natbib}
\usepackage[colorlinks=true,linkcolor=blue!50!black,citecolor=green!40!black,urlcolor=blue!60!black]{hyperref}

\setlist{nosep,leftmargin=*}
\captionsetup{font=small,labelfont=bf}
\renewcommand{\arraystretch}{1.12}
\setlength{\emergencystretch}{3em}
\newcolumntype{L}{>{\raggedright\arraybackslash}X}
\newcommand{\pending}[1]{\par\smallskip\noindent{\color{orange!70!black}\textbf{Results pending.} #1}\par\smallskip}
\newcommand{\nl}{\textbackslash n}
\newcommand{\p}[1]{``\texttt{#1}''}
\newcommand{\zbar}{\bar z}

\title{Small Attention-Head Sets That Carry Simple Tasks\\[3pt]
\large Choosing heads by what they decode, where they attend and what they cause}
\author{Filippo Merlo}
\date{17 November 2026 --- study design, results pending}

\begin{document}
\maketitle

\begin{abstract}
A language model that is told a task in words, or shown it by examples, writes information about the task into the outputs of some of its attention heads. Averaged over prompts and written into a prompt that does not state the task, the outputs of a few heads can make the model perform it. Which heads to choose is open: one recent method takes heads whose outputs decode to words that describe the task, another ranks heads by their causal effect on the answer. We compare these and new rules on 14 simple tasks in Llama-3.1-8B-Instruct, under one design in which every rule ranks all heads, sets of equal size are compared, and everything a rule uses is chosen on items, wordings and example formats that the test does not use. We first reproduce the two published rules in this design. We then change them one step at a time: heads are read through the layers that follow them with Head-JLens, task words are taken from what heads decode rather than from hand-written lists, heads are ranked by the attention they pay to each segment of the prompt, and the best combination of these features is selected without the test items. Finally, the selected rule and the published ones are tested on wordings of the task never used to choose heads, on prompts that give the task by examples, and with the set's outputs summed into one vector instead of written in head by head. The design separates what each change contributes and tests whether the chosen heads carry the task beyond the prompts they were chosen on.
\end{abstract}

% =====================================================================
\section{Introduction}\label{sec:intro}

Large language models can perform a task they are only told about in an instruction, or shown in a few input--output examples. Part of how they do so has been localised: when a prompt states a task, the outputs of a small number of attention heads at the last prompt token carry information about it. Averaged over many prompts of the task and added to a prompt that does not state it, these outputs make the model perform the task \citep{todd2024function,hendel2023incontext}. Such a \emph{function vector} is a compact handle on how a model represents a task. It can be read, compared between tasks and between ways of giving a task, and used to steer the model.

Which heads make up such a representation depends on how they are chosen. \citet{todd2024function} and \citet{davidson2025different} rank heads by their causal effect: how much setting one head to its mean output raises the probability of the correct answer on prompts from which the task information has been removed. This measures what matters for the answer, but it needs one patched forward pass per head and prompt, and it depends on how the task information is removed. \citet{yang2026shared} instead read each head's output in the vocabulary and keep the heads whose top tokens include words that describe the task, such as \emph{opposite} for antonyms or \emph{capital} for country--capital. This is cheap and readable, but it depends on a hand-written list of task words and on the readout. The readout they use, the logit lens applied to one head, sees only what a head writes directly into the output vocabulary and not what the layers after it make of that write.

The two studies also disagree on a question that the choice of heads bears on: whether a task given by instructions and the same task given by examples are carried by the same heads. \citet{davidson2025different} find that heads chosen by causal effect on instructions and on examples overlap little. \citet{yang2026shared} find that most lexical task heads are shared between the two prompting styles. The methods differ in more than the selection criterion, however: in the number of heads, in whether one set serves all tasks or each task has its own, in how the set is written into the prompt, and in which prompts the heads are chosen on and tested on. These differences make the results hard to compare.

We study the choice of heads under one design that every condition shares. Every rule ranks all heads, and sets of the same size are compared, from 10 to 120 heads, both as one set per task and as one set for all tasks. Everything a rule uses is chosen on one half of each task's items and tested on the other, and the tests beyond the main setting use wordings and example formats that were never used to choose heads. Within this design we proceed in three steps.
\begin{enumerate}
\item We reproduce the lexical task heads of \citet{yang2026shared} and the instruction function vectors of \citet{davidson2025different} (Sec.~\ref{sec:published}).
\item We change these rules one step at a time (Sec.~\ref{sec:rules}). We read heads with Head-JLens \citep{merlo2026headjlens}, which maps a head's output through the computation of the layers that follow it. We take the task words from what heads decode rather than from a hand-written list. We rank heads by the attention the last token pays to each segment of the prompt, following evidence that in-context learning gathers task information at template and function-word tokens \citep{bai2026performance}. We then choose among all combinations of these features by their accuracy on held-out items, without using the items they are reported on.
\item We test the selected rule and the published ones beyond the prompts they were chosen on (Sec.~\ref{sec:beyond}): on other wordings of each task, on prompts that give the task by examples, and with the set's outputs summed into a single vector as in a function vector.
\end{enumerate}

\paragraph{Questions.}
\begin{enumerate}[label=Q\arabic*.]
\item Do the two published selection rules reproduce when they are evaluated in the same way?
\item Which of our changes to them, taken one at a time, give sets that carry the task better at the same size?
\item To which segments of the prompt do the heads that carry a task attend, and which combination of decoding and attention picks the best sets?
\item Do the chosen sets carry the task when it is worded differently, when it is given by examples, and when the set is written in as one summed vector?
\end{enumerate}

% =====================================================================
\section{Background}\label{sec:background}

\subsection{Task information in attention-head outputs}\label{sec:bg-fv}

In a transformer, every layer adds its output to a residual stream that the next layers read. The output of an attention layer is a sum of one term per head \citep{elhage2021mathematical}: if $z_{l,h}$ is the output of head $h$ in layer $l$ before the output projection and $W_O^{(h)}$ is that head's block of the projection, the head writes
\begin{equation}
c_{l,h} = W_O^{(h)} z_{l,h}
\end{equation}
into the stream. Each head's contribution can therefore be read, replaced or removed on its own.

\citet{hendel2023incontext} show that in in-context learning the demonstrations are compressed into a single vector in the residual stream at an intermediate layer, which, patched into a prompt that holds only a query, makes the model apply the demonstrated task. \citet{todd2024function} locate such information in a small set of attention heads. They average each head's output at the last token over many prompts of a task. The average keeps what the prompts have in common, namely the task and the format, while what depends on each query largely cancels out. They then choose the heads whose average has the largest causal effect (Sec.~\ref{sec:bg-cie}), and sum their averaged writes into a function vector. Added to the residual stream at an early-to-middle layer, the function vector makes the model perform the task on zero-shot prompts and in natural text. \citet{yin2025which} find that such function-vector heads, more than induction heads, account for few-shot accuracy in larger models. \citet{yang2026localizing} separate heads that recognise which task a prompt shows from heads that learn its input--output mapping. In all of these accounts, a task is represented, at least in part, by the outputs of a small number of identifiable heads at the last prompt token.

\subsection{Stating a task and showing it}\label{sec:bg-styles}

A task can be given in two styles: by an instruction that describes it, or by examples that show it. Models do not learn equally well from both; on novel tasks they often learn more reliably from rules than from examples \citep{fu2026rules}. Whether the two styles produce the same internal representation is debated. \citet{davidson2025different} build function vectors from instructions as well as from demonstrations and find that the heads with the largest causal effect differ between the two styles: of the 20 top heads of each, only a few are shared, and those lie in later layers. \citet{yang2026shared} find instead that heads whose outputs decode to words describing the task, which they call \emph{lexical task heads}, are mostly shared between instructions and examples, and largely disjoint from function-vector heads. \citet{opielka2026causality} report that vectors chosen for their causal effect are not invariant to the format in which a task is given. A related observation is that a model's accuracy on a task changes with the wording of the instruction \citep{zhu2024promptrobust}, though part of the measured change can come from rigid answer matching rather than from the model \citep{hua2025flaw}. A set of heads that represents a task should therefore be tested on wordings and formats other than those it was chosen on.

\subsection{Reading a head in the vocabulary}\label{sec:bg-lens}

The \emph{logit lens} reads an intermediate state of the residual stream by applying the model's final normalisation and unembedding to it, as if no further layers followed \citep{nostalgebraist2020logitlens}. Applied to one head, it scores each token $v$ of the vocabulary by
\begin{equation}
s^{\text{logit}}_{l,h}(v) = q_v^\top c_{l,h},
\end{equation}
where $q_v$ is the unembedding row of $v$ scaled by the final normalisation's gain. This is the per-head readout of \citet{yang2026shared}, which we call \emph{Logit-head}. It measures only the head's direct effect on the output. A head in a middle layer may write a direction that later layers turn into a word, and the logit lens does not see this.

The Jacobian lens of \citet{gurnee2026workspace} reads a state through its first-order effect on the model's final layer. A small change to an intermediate state passes through the remaining layers and changes the final residual stream; to first order this dependence is a Jacobian. Averaged over many texts and positions, the Jacobian becomes one fixed map per layer, and an intermediate state is read by applying that map before the unembedding. Head-JLens \citep{merlo2026headjlens} places this construction at the concatenated head outputs of a layer, before they are summed, so that each head is read through its own columns $J_{l,h}$ of the fitted map:
\begin{equation}
s^{\text{J}}_{l,h}(v) = q_v^\top J_{l,h}\, z_{l,h}.
\end{equation}
We call this readout \emph{J-head}. It reads a head through the computation that follows it, and at most layers it ranks intermediate words that a prompt never states higher than Logit-head does \citep[Sec.~3]{merlo2026headjlens}. Under either readout, a head \emph{decodes} a word on a prompt when one of the word's token forms is among its ten highest-scoring tokens.

\subsection{Measuring a head's causal role}\label{sec:bg-cie}

The causal indirect effect of a head \citep{todd2024function} measures how much the head's typical output on prompts that give the task restores the correct answer on prompts that do not. Let $\tilde p$ be a prompt from which the task information has been removed, $y$ the first token of the correct answer, and $\zbar_{l,h}$ the head's mean output over prompts that give the task. Then
\begin{equation}
\mathrm{CIE}_{l,h} = \mathbb{E}_{\tilde p}\Big[\, P\big(y \mid \tilde p;\ z_{l,h} \leftarrow \zbar_{l,h}\big) - P\big(y \mid \tilde p\big) \Big].
\end{equation}
For examples, \citet{todd2024function} remove the task by shuffling the answers among the demonstrations. For instructions, \citet{davidson2025different} replace the instruction by uninformative text of the same length and log-probability: token sequences sampled to be as probable as the instruction, passages of ordinary text, and instructions of other tasks. The score needs one patched forward pass for every head and prompt.

\subsection{Where the last token attends}\label{sec:bg-attention}

A head's output at the last token is a weighted sum of what it reads from earlier positions, with the weights given by its attention row. Attention therefore tells where a head takes its information from. \citet{bai2026performance} find that in in-context learning the tokens that matter most for performance are template tokens and stopwords rather than content words, and that the model gathers task information at them. Many heads also put most of their attention on the first token when they have nothing specific to read \citep{xiao2024efficient}. Attention to the parts of a prompt that state the task, and little attention to the first token, are thus cheap signals that a head reads the task from the prompt. Unlike a causal score, they need only one forward pass per prompt.

% =====================================================================
\section{Experimental setup}\label{sec:setup}

\subsection{Model and tasks}\label{sec:model}

We use Llama-3.1-8B-Instruct in bfloat16, the main model of \citet{yang2026shared}. It has $L=32$ layers with $H=32$ heads each, 1{,}024 heads in all. Its Head-JLens maps were fitted on 1{,}000 WikiText-103 passages \citep{merlo2026headjlens}.

The tasks are the 14 simple tasks shared by \citet{yang2026shared} and \citet{davidson2025different}, with the datasets of \citet{todd2024function}. We group them by type:
\begin{itemize}
\item \textbf{word relations}: antonym, synonym;
\item \textbf{morphology}: singular--plural, present--past;
\item \textbf{facts}: country--capital, country--currency, landmark--country, park--country, product--company;
\item \textbf{sequences}: next-item, prev-item;
\item \textbf{translations}: English--French, English--German, English--Spanish.
\end{itemize}
The datasets hold between 197 (country--capital) and about 5{,}200 items (English--Spanish).

\subsection{Prompts}\label{sec:prompts}

All prompts are raw text without the chat template, and decoding is greedy. There are three families of task-stating prompts and one test prompt.
\begin{itemize}
\item \textbf{Five instruction templates} per task, those of \citet{yang2026shared}: one cloze template (\p{The antonym of \{x\} is}) and four of the form \p{<instruction> Q: \{x\}\nl A:}. These are used to choose heads and to compute means in Sections~\ref{sec:published} and~\ref{sec:rules}.
\item \textbf{Thirty further wordings} per task, written by GPT-5.5 and checked by script, in six groups of five: three forms (cloze, question--answer, label) crossed with whether the wording names the task, that is, contains a word of \citeauthor{yang2026shared}'s term list or the task's name (App.~\ref{app:prompts}). They are split once into a choosing part P and a test part Q of 15 wordings each, with three and two of every group in turn.
\item \textbf{Six example formats} at five shots (Table~\ref{tab:formats}), likewise split into three choosing formats P and three test formats Q, so that each part holds one labelled format, one arrow-like format and one with a plain separator. Demonstrations are drawn at random among the other items of the same half.
\item \textbf{The test prompt} is the zero-shot prompt \p{ x:}. It contains the query and nothing that states the task.
\end{itemize}

\begin{table}[tbp]
\centering\small
\begin{tabularx}{\textwidth}{@{}cLLc@{}}
\toprule
 & Format & Prompt (two of five shots shown) & Part \\
\midrule
F1 & colon and semicolon \citep{yang2026shared} & \p{ hot: cold; tall: short; wet:} & P \\
F2 & question--answer \citep{todd2024function} & \p{Q: hot\nl A: cold\nl\nl Q: tall\nl A: short\nl\nl Q: wet\nl A:} & P \\
F3 & arrow & \p{hot -> cold\nl tall -> short\nl wet ->} & P \\
F4 & input--output labels & \p{Input: hot\nl Output: cold\nl\nl Input: wet\nl Output:} & Q \\
F5 & equals & \p{hot = cold, tall = short, wet =} & Q \\
F6 & bar & \p{hot | cold\nl tall | short\nl wet |} & Q \\
\bottomrule
\end{tabularx}
\caption{Example formats, shown on antonym. P formats are used to choose heads, Q formats only to compute the means of the test.}
\label{tab:formats}
\end{table}

\subsection{Held-out protocol}\label{sec:heldout}

\paragraph{Items.} Each task's items are shuffled once and cut into two halves, $A$ and $B$. In fold~1 everything that depends on items is computed on $A$ and the sets are tested on $B$; fold~2 is the reverse. Reported accuracies are means over the two folds. ``Everything that depends on items'' covers the prompts whose heads are read, the word lists, attention, causal scores, generic heads, and the task and control means.

\paragraph{Prompts.} Heads are never chosen on Q wordings or Q formats; these only supply the means written into the test prompt. Every choosing pool uses up to 300 correctly answered prompts per task, spread evenly over its templates or formats, so that instructions and examples are chosen on the same amount of material. The causal score keeps the counts of \citet{davidson2025different}: 100 prompts for the means and 25 for each kind of uninformative text.

\paragraph{Cross-fitting.} Some choices are made by accuracy: the scaling $\alpha$ and, in Sec.~\ref{sec:select}, the combination of features. Choosing the best of several candidates on the same items it is then reported on gives an optimistic estimate, because the winner is partly the candidate that happened to do well on those items. We therefore cross-fit. A choice for fold~1 is made on fold~2's accuracy and the reverse. Fold~2 is tested on half $A$, which fold~1 uses only for choosing, so no outcome on the items a fold is scored on enters its choices.

\subsection{Writing a head set into a prompt}\label{sec:patching}

\paragraph{Means.} For a task $t$, a template or format $w$ and a head $(l,h)$, the mean $\zbar^{\,t,w}_{l,h}$ is the head's output at the last token averaged over the first 100 correctly answered prompts of $w$ on the choosing half.

\paragraph{Head by head.} Following \citet[Sec.~2.3.2]{yang2026shared}, every head of a set $S$ has its output at the last token of the test prompt replaced by a scaled mean,
\begin{equation}
z_{l,h} \leftarrow \alpha\, \zbar^{\,t,w}_{l,h} \qquad \text{for all } (l,h)\in S,
\end{equation}
at its own layer; the layers after it compute on the changed stream, and the reply is generated from the patched prompt. This is the default intervention.

\paragraph{Summed.} Following \citet{davidson2025different}, the means are taken through their output projections, summed, and added to the residual stream $x_{\ell^*}$ after layer $\ell^* \approx L/3$ at the last token:
\begin{equation}
x_{\ell^*} \leftarrow x_{\ell^*} + \alpha \sum_{(l,h)\in S} W_O^{(h)}\, \zbar^{\,t,w}_{l,h}.
\end{equation}
The two interventions differ in two ways. Head by head replaces a head's output, whereas the summed vector adds to the stream. And head by head writes each head at its own layer, whereas the summed vector writes all heads at one layer, earlier than heads above $\ell^*$ and later than heads below it. The summed vector therefore carries the direction that the set writes, but not where in the network each part is written. On our model $\ell^*=11$.

\paragraph{Scaling.} $\alpha\in\{1,2\}$, cross-fitted (Sec.~\ref{sec:heldout}): each fold uses the value that was better on the other fold for the same task and set.

\subsection{Evaluation}\label{sec:evaluation}

\paragraph{Test.} The patched prompts are the test prompts \p{ x:} of the items the model answers correctly with one example in format F1, up to 100 per task and fold. Without a patch these prompts are almost never answered correctly, since they do not state the task, so accuracy measures how much task information the set brings in. Each item is patched with every mean of the test templates or formats in turn, and accuracy is averaged over them. Replies stop at a newline, the example delimiter or 16 tokens.

\paragraph{Scoring.} A reply is correct when it begins with the whole reference answer as words, ignoring case, accents, punctuation and a leading article. Matching only the first token, as \citet{yang2026shared} do, would count any number as correct on next-item and prev-item and a capitalised German noun as wrong.

\paragraph{Controls.} Two controls separate task information from other effects of patching.
\begin{itemize}
\item The \emph{control mean}: the same heads patched with the means of a related task on the same templates or formats, using the pairing of \citet{yang2026shared} (e.g.\ antonym with synonym; App.~\ref{app:prompts}). A set carries its task when its own mean raises accuracy and the control mean does not.
\item \emph{Random sets}: for every rule, scope and size, three sets with as many heads in each layer as the rule's set, drawn uniformly among the heads of those layers and scored at the rule's $\alpha$. They test whether the choice of heads matters beyond their layers. Their overlap with the rule's set is reported.
\end{itemize}

\subsection{Rankings, set sizes and scopes}\label{sec:scopes}

Every rule assigns each head a score on each task, from which ranks $\rho_t(h)$ are formed (1 is best; ties share their mean rank). The set is the top $N$ heads, $N\in\{10,20,40,60,80,100,120\}$, with remaining ties broken by the lower layer, then head. Sets are compared only at equal $N$. A rule that cannot fill $N$ heads at some size is not run there.

\begin{itemize}
\item \textbf{Per task}: each task gets the top $N$ of its own ranking.
\item \textbf{Shared}: one set serves all tasks. Each head's ranks are averaged over the tasks, $\bar\rho(h) = \frac{1}{|T|}\sum_{t} \rho_t(h)$, and the top $N$ by $\bar\rho$ form the set. Each task is still patched with its own means. Averaging ranks rather than scores keeps tasks with larger scores from dominating.
\end{itemize}
Accuracy is reported as the mean over tasks and by task type. Every condition is one greedy run, and we report accuracies only.

% =====================================================================
\section{Two published ways of choosing heads}\label{sec:published}

We reproduce both methods in two forms: their own form, which fixes the set as the source study does, and a ranking that is compared with the other rules at every size and in both scopes. Both are chosen on the five instruction templates and patched with their means.

\subsection{Lexical task heads}\label{sec:yang}

\paragraph{Own form.} \citet{yang2026shared} write a list of task-describing terms for each task, with their token variants. A head is a lexical task head for a task if, on at least 10\% of the first 100 correctly answered prompts of the task's most accurate template, at least one token of the list is among its ten highest Logit-head tokens. The set is per task, and its size varies with the task.

\paragraph{Ranking.} \textbf{Yang terms} ranks heads by the same share of prompts. Heads that never decode a term are not ranked, so at large $N$ this rule may not fill the set.

\subsection{Instruction function vectors}\label{sec:davidson}

\paragraph{Own form.} \citet{davidson2025different} score every head by its causal indirect effect on instruction prompts (Sec.~\ref{sec:bg-cie}), with the instruction replaced by the three kinds of uninformative text, averaged over kinds and over tasks. The 20 heads with the largest average form one set for all tasks, written in as one summed vector with $\alpha=1$. We compute the score on the four question--answer templates, whose instruction can be replaced while the query is kept; the cloze template has no instruction apart from the query.

\paragraph{Ranking.} \textbf{Causal} ranks heads by the same score, per task or averaged over tasks by rank (Sec.~\ref{sec:scopes}), and is patched head by head like every other rule.

\paragraph{Settings of the source studies.} Each method is also run in the settings of its source, which differ from ours in padding, answer scoring, scaling and the patched prompts (App.~\ref{app:fidelity}).

\pending{Own-form accuracies of both methods; Yang terms and Causal at every size and in both scopes; the settings of the source studies.}

% =====================================================================
\section{Building selection rules from readouts and attention}\label{sec:rules}

Starting from the published rules, we change one thing at a time (Table~\ref{tab:ladder}). Each step keeps everything else of the step before, so a change in accuracy at equal size can be attributed to that step. The rules of the ladder are run at every size and in both scopes, with their control means and random sets.

\begin{table}[tbp]
\centering\small
\begin{tabularx}{\textwidth}{@{}clLl@{}}
\toprule
 & Change & Reason & Rule \\
\midrule
1 & Logit-head $\to$ J-head & reads a head through the layers that follow it (Sec.~\ref{sec:bg-lens}) & Yang terms (J) \\
2 & hand-written terms $\to$ words from the data & removes the dependence on a list written in advance & D \\
3 & attention to each prompt segment & where a head reads from (Sec.~\ref{sec:bg-attention}) & one rule per segment \\
4 & best combination of features & combines decoding and attention, chosen on held-out accuracy & Selected \\
5 & generic heads removed from the causal ranking & separates task-specific heads from heads that help every task alike & Causal $-$ generic \\
\bottomrule
\end{tabularx}
\caption{The steps from the published rules to ours. Steps 1--4 start from lexical task heads, step 5 from the causal ranking.}
\label{tab:ladder}
\end{table}

\subsection{Reading heads through the layers that follow}\label{sec:jhead}

\textbf{Yang terms (J)} is Yang terms with J-head in place of Logit-head. A head in a middle layer whose write becomes a task word only after later layers transform it is found by J-head and missed by Logit-head (Sec.~\ref{sec:bg-lens}).

\subsection{Task words from the data}\label{sec:lists}

A hand-written list fixes in advance which words count as describing a task. We instead take each task's words from what its heads decode, on the choosing half alone:
\begin{enumerate}
\item every head's ten highest J-head tokens at the last token of the choosing prompts;
\item a token counts as a word if, trimmed and lower-cased, it consists of letters or digits only, is not one of 18 common stop words, and has at least three characters when it contains a Latin letter;
\item a word's \emph{share} on a task is the fraction of (prompt, head) pairs whose top ten contain it;
\item the 50 words with the largest share on 1{,}000 WikiText-103 passages, read in the same way, are removed as \emph{generic words}: heads decode them on any text;
\item the task's list is its 30 remaining words with the largest share.
\end{enumerate}
\textbf{D} ranks heads by the share of choosing prompts on which they decode a word of the list.

\subsection{Attention to the segments of the prompt}\label{sec:segments}

We split every choosing prompt into segments, so that each token belongs to exactly one, and rank heads by how much attention the last token pays to each. On instruction prompts the segments are assigned by script, in this order of priority:
\begin{enumerate}
\item the first token;
\item \emph{task words}: instruction tokens, outside the query, that match a word of the task's list;
\item the \emph{query} $x$;
\item the \emph{last token}, where the head is read;
\item \emph{markers}: the scaffold around the query (\p{ Q:} before it and \p{\nl A} after it in question--answer templates, the labels in label templates; cloze templates have none);
\item \emph{punctuation}: tokens without a letter or digit;
\item \emph{stopwords}: NLTK's English stopword list;
\item \emph{other words}: every remaining instruction token.
\end{enumerate}
For head $(l,h)$ and segment $s$, the score is the attention mass that the head's last-token row puts on the segment,
\begin{equation}
m_{l,h}(s) = \mathbb{E}_{p}\Big[\sum_{j\in s(p)} a_{l,h}(p)_{n,j}\Big],
\end{equation}
where $a_{l,h}(p)_{n,j}$ is the attention from the last position $n$ to position $j$ and the mean runs over the choosing prompts in which the segment is not empty. Instead of attention to the first token, which marks a head with nothing specific to read (Sec.~\ref{sec:bg-attention}), we use its complement, the attention to every other token (\emph{non-first}). All features rank heads by high values. With D, there are nine features (Table~\ref{tab:features}).

\begin{table}[tbp]
\centering\small
\begin{tabularx}{\textwidth}{@{}lLL@{}}
\toprule
Feature & Instruction prompts & Example prompts (Sec.~\ref{sec:examples}) \\
\midrule
D & decodes a word of the task's list & same \\
non-first & attention to every token but the first & same \\
task words & instruction tokens matching the list & -- \\
other words & remaining content words of the instruction & -- \\
stopwords & stopwords of the instruction & -- \\
punctuation & tokens without a letter or digit & -- \\
markers & scaffold around the query & -- \\
inputs & -- & the demonstrations' inputs \\
outputs & -- & the demonstrations' outputs \\
separators & -- & tokens between and around demonstrations (\p{:}, \p{;}, \p{Q:}, \p{A:}, \p{->}, \p{=}, \p{|}, labels, newlines) \\
query & the query $x$ & same \\
last token & the read-out position & same \\
\midrule
Features & 9 & 7 \\
Combinations & 511, of which 255 without D & 127, of which 63 without D \\
\bottomrule
\end{tabularx}
\caption{Features from which the selection rules are built.}
\label{tab:features}
\end{table}

\subsection{Choosing a combination without the test items}\label{sec:select}

A combination $F$ of features ranks heads by the mean of its features' within-task ranks, $\rho^F_t(h) = \frac{1}{|F|}\sum_{f\in F}\rho^f_t(h)$. For a shared set, each feature's ranks are first averaged over tasks and ranked again, and the combination averages these. Every non-empty subset of the nine features is a candidate: D alone, each segment alone, every combination of segments without D, and every combination with D, 511 per scope.

Each candidate is built at 40 heads and patched in both folds with the task means. Candidates that give the same 40 heads are patched once. The control mean is not used for choosing and is run only for the selected rule. For fold~1, \textbf{Selected} is the candidate with the highest accuracy on fold~2, at the better of $\alpha=1$ and $2$; for fold~2 the reverse; separately per task and shared. The two folds may select different combinations, and both are reported. Selected is then run at every size.

The same runs show what each feature contributes: its accuracy alone at 40 heads, and the change in accuracy when it is added to each combination that lacks it.

\subsection{Generic heads in the causal ranking}\label{sec:generic}

A causal score averaged over tasks rewards heads that help on every task in the same way, for instance by signalling that a short answer is to follow. The mean outputs of such heads are much the same for every task, so they add a direction that does not depend on the task. We call them \emph{generic heads} and identify them by how similar their readings are across tasks. A head's generic score is
\begin{equation}
g_{l,h} = \frac{1}{\binom{|T|}{2}} \sum_{t<t'} \cos\big(n_{t}(l,h),\, n_{t'}(l,h)\big),
\end{equation}
where $n_t(l,h)$ counts how often each token is among the head's ten highest J-head tokens on task $t$'s choosing prompts. Among the 30 top heads of the shared causal ranking, the sorted generic scores are split into a high and a low group with the least spread within each group. The high group are the generic heads, unless it holds more than half of the 30. \textbf{Causal $-$ generic} removes them from the causal ranking before the top $N$ is taken. The step tests whether the causal ranking carries task-specific information once heads that help every task alike are set aside.

\pending{The steps of Table~\ref{tab:ladder} at every size and in both scopes, by task type; the selected combinations; each feature alone and added; overlap of Selected with Causal; layers of the chosen heads.}

% =====================================================================
\section{Beyond the prompts the heads were chosen on}\label{sec:beyond}

The rules of Sec.~\ref{sec:rules} are chosen and tested on the same five instruction templates. We test five of them further: Yang terms (J), D, Selected, Causal and Causal $-$ generic, with their control means and random sets, at every size and in both scopes. The test prompt stays \p{ x:}; each test changes one thing with respect to Sec.~\ref{sec:rules}: the prompts the means come from, the prompts the heads are chosen on, or the intervention (Table~\ref{tab:scenarios}).

\begin{table}[tbp]
\centering\small
\begin{tabularx}{\textwidth}{@{}lLLl@{}}
\toprule
Test & Heads chosen on & Means from & Changed \\
\midrule
Sec.~\ref{sec:rules} & five instruction templates & five instruction templates & -- \\
Wordings, transfer & five instruction templates & 15 Q wordings & means \\
Wordings, re-chosen & five templates + 15 P wordings & 15 Q wordings & choosing prompts \\
Examples, transfer & five instruction templates & 3 Q formats & prompting style of the means \\
Examples, re-chosen & 3 P formats & 3 Q formats & prompting style of the choice \\
Summed vector & as above & as above & intervention \\
\bottomrule
\end{tabularx}
\caption{Tests beyond the choosing prompts.}
\label{tab:scenarios}
\end{table}

\subsection{Unseen wordings}\label{sec:wordings}

\paragraph{Transfer.} The sets of Sec.~\ref{sec:rules}, unchanged, are patched with the means of the 15 Q wordings, which played no part in choosing any head, list or scale. If a set carries the task rather than the phrasing of the templates it was chosen on, its accuracy should hold. Accuracy is also reported by wording group, since forms that do not name the task, or that differ from the question--answer layout, may be harder.

\paragraph{Re-chosen.} The same rules are run again on the five templates together with the 15 P wordings: the word lists, attention, causal scores on the question--answer wordings, generic heads and the combination screen are all recomputed. The sets are tested with the same Q means. Comparing them with the transferred sets at equal size shows whether heads chosen on more varied wordings carry unseen wordings better. The share of heads the two sets have in common is reported.

\subsection{Tasks given by examples}\label{sec:examples}

\paragraph{Transfer.} The sets of Sec.~\ref{sec:rules} are patched with the means of the three Q example formats. This tests the claim of \citet{yang2026shared} that heads that carry a task stated in words also carry it when it is shown by examples.

\paragraph{Re-chosen.} Every rule is rebuilt on five-shot prompts of the three P formats. The segments of an example prompt differ from those of an instruction (Table~\ref{tab:features}): the first token, the demonstrations' inputs, their outputs, the separators between and around them, the query and the last token. Task words are not used, since demonstrations seldom contain words that describe the task. The combination screen is run over the seven example features as in Sec.~\ref{sec:select}. The causal score is the one \citet{davidson2025different} use for demonstrations, that of \citet{todd2024function}: the effect of a head when the answers are shuffled among the demonstrations. The sets are tested with the same Q means as the transferred ones.

\paragraph{Overlap.} For each rule, scope and size, the share of heads that the set chosen on instructions and the set chosen on examples have in common, $|S_I\cap S_E|/N$, is compared with the share expected for two random sets with the same number of heads in each layer, $\frac{1}{N}\sum_l n_l m_l / H$, where $n_l$ and $m_l$ are the two sets' numbers of heads in layer $l$. \citet{davidson2025different} find little overlap for the causal score and \citet{yang2026shared} much for lexical task heads; with all rules measured in the same way, the comparison shows whether the difference lies in the rule or in the setting.

\subsection{One summed vector or head by head}\label{sec:summed}

The sets of Sec.~\ref{sec:rules}, with the means of the five templates, and the sets re-chosen on example prompts, with the Q format means, are written in as one summed vector (Sec.~\ref{sec:patching}) and compared with head-by-head patching at the same sizes, means and scaling. If the task were carried only by the direction a set writes, the two interventions would give similar accuracy; if it depends on where each head writes, head by head should be better, and increasingly so for larger sets, which span more layers.

\pending{Each test at every size and in both scopes, with control means and random sets; accuracy by wording group and by task type; overlap between sets chosen on instructions and on examples.}

% =====================================================================
\section{Limitations}\label{sec:limitations}

\paragraph{Scope.} One model and 14 simple tasks whose answers are single words or short phrases. Every condition is one greedy run, and accuracies are reported without intervals, so differences of a point or two between sets are not established.

\paragraph{Scaling.} Only $\alpha\in\{1,2\}$ is compared, which may underrate sets whose best scaling is higher, in particular summed vectors.

\paragraph{Prompt material.} The 30 further wordings were generated by a language model and checked by script, not by hand. Segments are assigned by script, and the task-word segment depends on the task's list, so it is empty in templates that use none of its words. Example prompts have one format per mean, and their separators differ between formats.

\paragraph{Selection.} Selected is the best of many candidates on the other fold. Cross-fitting keeps the test items out of the choice, but with two folds the choice rests on half of the items, and the folds may select different combinations.

\paragraph{Generic heads.} They are defined among the top heads of the shared causal ranking and removed only from it; whether heads of the same kind harm other rules is not tested.

% =====================================================================
\section*{Data and code availability}

The task datasets are those of \citet{todd2024function} as released by \citet{yang2026shared}; the Head-JLens maps are those of \citet{merlo2026headjlens}. Code, sets and every patched reply with its expected answer are kept in the project repository.

\bibliographystyle{plainnat}
\bibliography{references}

% =====================================================================
\clearpage
\appendix

\section{Settings of the source studies}\label{app:fidelity}

\begin{table}[H]
\centering\small
\begin{tabularx}{\textwidth}{@{}>{\raggedright\arraybackslash}p{3.2cm}LL@{}}
\toprule
Setting & \citet{yang2026shared} & This study \\
\midrule
Model & 13 models, main text Llama-3.1-8B-Instruct & Llama-3.1-8B-Instruct, bfloat16 \\
Tasks & 17 simple tasks & the 14 shared with \citet{davidson2025different} \\
Instruction templates & five per task & the same five; 30 further wordings (Sec.~\ref{sec:prompts}) \\
Batching & left-padded batches, padding attended & padding masked \\
Readout & Logit-head & Logit-head and J-head \\
Head criterion & $\geq$1 term in the top 10 on $\geq$10\% of 100 correct prompts & same (own form); share of prompts (ranking) \\
Answer score & first generated token & whole answer as words \\
Scaling & 15 values from $-10$ to $10$ & 1 and 2, cross-fitted \\
Patched prompts & zero-shot prompts (paper); wrongly answered 1-shot prompts (released code) & zero-shot prompts \\
\bottomrule
\end{tabularx}
\caption{Lexical task heads: the source study's settings and ours.}
\label{tab:fid-yang}
\end{table}

\begin{table}[H]
\centering\small
\begin{tabularx}{\textwidth}{@{}>{\raggedright\arraybackslash}p{3.2cm}LL@{}}
\toprule
Setting & \citet{davidson2025different} & This study \\
\midrule
Model & 12 models incl.\ Llama-3.1-8B-Instruct & Llama-3.1-8B-Instruct \\
Tasks & 50 datasets of \citet{todd2024function} & the 14 shared with \citet{yang2026shared} \\
Instructions & five best of about 200 generated instructions, short and long & the four question--answer templates of \citet{yang2026shared} \\
Mean outputs & 100 correct prompts, 20 per instruction & 100 correct prompts, 25 per template \\
Uninformative text & three kinds, matched in length and log-probability & same \\
Heads & top 20, one set for all tasks & top 20 (own form); 10 to 120, per task and shared (ranking) \\
Intervention & summed vector after layer $L/3$, weight 1 & same (own form); head by head and summed, $\alpha$ 1 and 2 \\
Evaluation & zero-shot prompt with the query & \p{ x:}, whole-answer score, control mean \\
\bottomrule
\end{tabularx}
\caption{Instruction function vectors: the source study's settings and ours.}
\label{tab:fid-dav}
\end{table}

\section{Prompt material}\label{app:prompts}

\begin{table}[H]
\centering\small
\begin{tabularx}{\textwidth}{@{}lL@{}}
\toprule
Group & Example (antonym) \\
\midrule
\multicolumn{2}{@{}l}{\textit{Instruction templates of \citet{yang2026shared}}} \\
cloze & \p{The antonym of \{x\} is} \\
question--answer & \p{Output the antonym of the word in the question. Q: \{x\}\nl A:} \\
\midrule
\multicolumn{2}{@{}l}{\textit{Further wordings (five per group)}} \\
cloze, names the task & \p{An antonym for \{x\} would be} \\
cloze, does not name it & \p{The contrasting word for \{x\} is} \\
question--answer, names the task & \p{Provide the antonym for the given word. Q: \{x\}\nl A:} \\
question--answer, does not name it & \p{Provide a word with a contrary meaning to the given term. Q: \{x\}\nl A:} \\
label, names the task & \p{Word: \{x\}\nl Antonym:} \\
label, does not name it & \p{Word: \{x\}\nl Contrary word:} \\
\bottomrule
\end{tabularx}
\caption{Instruction templates and further wordings, one per group.}
\label{tab:wordings}
\end{table}

\paragraph{Control pairs} \citep{yang2026shared}: antonym and synonym with each other; singular--plural and present--past with each other; next-item and prev-item with each other; English--French and English--German with each other, English--Spanish with English--German; country--capital, country--currency and product--company with park--country; park--country with country--currency; landmark--country with product--company.

\end{document}
